Si presenta dunque l'architettura di \emph{BlueAgent}, una piattaforma
di Cyber Range in esecuzione come servizio \emph{HTTP} che espone
\emph{API REST} per il deployment, la gestione del ciclo di vita e la
verifica automatizzata delle challenge, a partire da definizioni
dichiarative (\emph{Infrastructure as Code}) predisposte dagli
organizzatori della gara. Tali API REST possono essere consumate da un
qualsiasi applicativo esterno denominato CTF Manager, come ad esempio un
plugin per la piattaforma CTFd\footnote{\label{fn:ref_13}CTFd LLC, \emph{CTFd: Open
  Source Capture The Flag Framework}, https://ctfd.io/}.

Il Cyber Range implementato da BlueAgent è pensato per scenari di Patch
Management da eseguire su un singolo host. Si propone di offrire un
ambiente basato su container di sistema, che garantisca un meccanismo di
verifica completamente inaccessibile ai concorrenti. Tale meccanismo si
fonda sulla clonazione all'occorrenza dell'intera infrastruttura del
team in una rete logica separata, sulla quale vengono eseguiti i comandi
di verifica definiti dall'oracolo.

I membri di ciascun team possono connettersi alla rete dei container a
loro assegnata mediante una VPN WireGuard, dove a ciascun partecipante
corrisponde un peer dedicato (ossia, coppia di chiavi pubblica e
privata).

È opportuno evidenziare che:

\begin{enumerate}
\def\labelenumi{\arabic{enumi}.}
\tightlist
\item
  Gli indirizzi IP dei container all'interno della rete di produzione di
  ciascun team sono statici e noti a priori (il che si applica alla
  maggior parte delle challenge).
\item
  La duplicazione istantanea sfrutta le proprietà Copy-on-Write (CoW)
  del filesystem ZFS. Inoltre viene impiegata una rete SDN
  (Software-Defined Network) per orchestrare router e switch logici
  dedicati all'ambiente di verifica.
\item
  BlueAgent non mantiene lo stato della competizione (quali i punteggi o
  la classifica dei team), delegando tale responsabilità al CTF Manager.
  BlueAgent opera esclusivamente come un orchestratore specializzato che
  traduce le specifiche delle challenge in container da eseguire.
\end{enumerate}

A scopo d'esempio, mostriamo ad alto livello come si crea una challenge:

\begin{enumerate}
\def\labelenumi{\arabic{enumi}.}
\tightlist
\item
  Il creatore della challenge definisce tramite un file YAML i
  container.
\item
  Il creatore della challenge, inserisce la nuova challenge tenendo in
  memoria tutte le informazioni necessarie a BlueAgent (inclusi i file
  dichiarativi)
\item
  Il giocatore di un determinato team accede alla challenge e tramite il
  CTF Manager richiede la creazione dei container.
\item
  BlueAgent si occupa di istanziare i container per il team.
\end{enumerate}

BlueAgent è basato su \emph{FastAPI}, per gestire in modo semplice ed
efficiente le richieste HTTP.

L'host su cui viene eseguito BlueAgent deve mettere a disposizione
dell'architettura le seguenti componenti:

\begin{itemize}
\tightlist
\item
  \textbf{Debian 13}: o una qualsiasi distribuzione Linux.
\item
  \passthrough{\lstinline!nftables!}: per filtrare i pacchetti nel
  kernel
\item
  \textbf{Python 3}: per eseguire BlueAgent.
\item
  \textbf{Pool ZFS}: filesystem che supporta il meccanismo
  \passthrough{\lstinline!COW!}.
\item
  \textbf{OVN ed Open vSwitch}: per gestire la rete di verifica dei
  container.
\item
  \textbf{Distrobuilder}: per generare le immagini dei container in modo
  dichiarativo.
\item
  \textbf{WireGuard}: per fornire una VPN come punto d'accesso ai
  partecipanti
\item
  \textbf{Incus}: per la gestione dei container.
\end{itemize}

Per comprendere passo-passo cosa permette di fare il sistema e come lo
fa sono elencati gli attori, gli use cases ed i requisiti di sistema.
